% Part: many-valued-logic
% Chapter: three-valued-logics
% Section: multiple-designation

\documentclass[../../../include/open-logic-section]{subfiles}

\begin{document}

\olfileid{mvl}{thr}{mul}

\olsection{指定值不止是 $\True$}

目前为止我们见过的逻辑，其指定真值集合都是 $V^+ = \{\True\}$，
也就是说，某公式算作真，当且仅当其真值为~$\True$。
但我们也可以认为，只要某公式不是~$\False$，它就算作真。
于是，例如通过在矩阵中规定 $V^+ = \{\True, \Undef\}$，就可以得到一种逻辑。

\begin{defn}
\emph{悖论逻辑}~$\LogLP$ 由以下矩阵定义：
\begin{enumerate}
  \item 标准命题语言 $\Lang L_0$，包含联结词
  $\lnot$、$\land$、$\lor$、$\lif$。
  \item 真值集合为 $V = \{\True, \Undef, \False\}$。
  \item $\True$ 与 $\Undef$ 是指定值，即 $V^+ = \{\True, \Undef\}$。
  \item 真值函数与强 Kleene 逻辑相同。
\end{enumerate}
\end{defn}

\begin{defn}
Halldén 的\emph{无意义逻辑}~$\LogHal$ 通过如下矩阵定义：
\begin{enumerate}
  \item 标准命题语言 $\Lang L_0$，包含联结词
  $\lnot$、$\land$、$\lor$、$\lif$ 以及一个一元联结词~$+$。
  \item 真值集合为 $V = \{\True, \Undef, \False\}$。
  \item $\True$ 与 $\Undef$ 是指定值，即 $V^+ = \{\True, \Undef\}$。
  \item 真值函数与弱 Kleene 逻辑相同，但额外增加一个“无意义”算子：
  \begin{center}
    \begin{tabular}{c|c}
    $\tf{+}$ & \\
    \hline
    $\True$ & $\False$ \\
    $\Undef$ & $\True$ \\
    $\False$ & $\False$
    \end{tabular}
  \end{center}
\end{enumerate}
\end{defn}

与它们共用相同真值表的 Kleene 逻辑不同，这些逻辑\emph{确实}有重言式。

\begin{prop}\ollabel{prop:LP-taut-CL}
 $\LogLP$ 的重言式与经典命题逻辑的重言式相同。
\end{prop}

\begin{proof}
  根据 \olref[syn][sub]{prop:mvl-cl}，若 $\Entails[\LogLP] !A$ 则
  $\Entails[\LogCL] !A$。为证其逆，我们证明：若存在赋值~$\pAssign v\colon \PVar \to \{\False, \True,
  \Undef\}$ 使得 $\pValue v(!A)[\LogKs] = \False$，则存在赋值~$\pAssign {v'}\colon \PVar \to \{\False, \True\}$ 使得 $\pValue {v'}(!A)[\LogCL] = \False$。由于 $\LogKs$ 和 $\LogLP$ 具有相同的特征真值函数，而 $\False$ 是 $\LogLP$ 中唯一非指定真值（这是 $\LogLP$ 与~$\LogKs$ 的唯一区别），这就为 $\LogLP$ 确立了结论。
  因此，如果 $\Entails/[\LogLP] !A$，则对某个赋值~$\pAssign v$，有 $\pValue v(!A)[\LogLP] = \pValue
  v[\LogKs](!A) = \False$。根据我们要证明的断言，$\pValue{v'}[\LogCL](!A) = \False$，即 $\Entails/[\LogCL] !A$。

  为证明此断言，我们首先如下定义 $\pAssign {v'}$：
  \[
    \pAssign {v'}(p) =
    \begin{cases}
    \True & \text{若 } \pAssign {v}(p) \in \{\True, \Undef\}\\
    \False & \text{否则}
  \end{cases}
  \]
  现在我们对 $!A$ 进行归纳，以证明 (a)~若 $\pValue v(!A)[\LogKs]
  = \False$ 则 $\pValue {v'}(!A)[\LogCL] = \False$，以及 (b)~若
  $\pValue v(!A)[\LogKs] = \True$ 则 $\pValue {v'}(!A)[\LogCL] =
  \True$。
  \begin{enumerate}
    \item 归纳基始：$!A \ident p$。根据
    \olref[syn][val]{defn:pValue}，$\pValue v(!A)[\LogKs]
    = \pAssign v(p) = \pValue {v'}(!A)[\LogCL]$，这蕴涵 (a) 和~(b)。

    对于归纳步骤，考虑以下情形：
    \item $!A \ident \lnot !B$。
    \begin{enumerate}
      \item 假设 $\pValue v(\lnot!B)[\LogKs] = \False$。根据
      $\tf{\lnot}[\LogKs]$ 的定义，$\pValue v(!B)[\LogKs]  =
      \True$。由归纳假设情况，得到 $\pValue
      {v'}(!B)[\LogCL]  = \True$，于是 $\pValue {v'}(\lnot !B)[\LogCL] = \False$。
      \item 假设 $\pValue v(\lnot!B)[\LogKs] = \True$。根据
      $\tf{\lnot}[\LogKs]$ 的定义，$\pValue v(!B)[\LogKs]  =
      \False$。由归纳假设情况，得到 $\pValue
      {v'}(!B)[\LogCL]  = \False$，于是 $\pValue {v'}(\lnot !B)[\LogCL] = \True$。
    \end{enumerate}

    \item $!A \ident (!B \land !C)$。
    \begin{enumerate}
      \item 假设 $\pValue v(!B \land !C)[\LogKs] = \False$。根据
      $\tf{\land}[\LogKs]$ 的定义，$\pValue v(!B)[\LogKs]  =
      \False$ 或 $\pValue v(!B)[\LogKs]  = \False$。由归纳假设情况，得到 $\pValue {v'}(!B)[\LogCL]  = \False$ 或 $\pValue {v'}(!C)[\LogCL]  = \False$，于是 $\pValue {v'}(!B \land
      !C)[\LogCL] = \False$。
      \item 假设 $\pValue v(!B \land !C)[\LogKs] = \True$。根据
      $\tf{\land}[\LogKs]$ 的定义，$\pValue v(!B)[\LogKs]  =
      \True$ 且 $\pValue v(!B)[\LogKs]  = \True$。由归纳假设情况，得到 $\pValue {v'}(!B)[\LogCL]  = \True$ 且 $\pValue {v'}(!C)[\LogCL]  = \True$，于是 $\pValue {v'}(!B \land
      !C)[\LogCL] = \True$。
    \end{enumerate}
  \end{enumerate}

  其余两种情况是类似的，留作练习。另一种方法是，上面的证明已对所有仅包含 $\lnot$ 和~$\land$ 的公式确立了结论。接下来可以利用以下事实：在 $\LogKs$ 和 $\LogCL$ 中，对任意 $\pAssign v$，$\pValue v(!B \lor
  !C) = \pValue v(\lnot(\lnot!B \land \lnot !C))$ 且 $\pValue v(!B \lif
  !C) = \pValue v(\lnot(!B \land \lnot !C))$。
\end{proof}

\begin{prob}
  完成 \olref[mvl][thr][mul]{prop:LP-taut-CL} 的证明，即对 $!A \ident (!B \lor !C)$ 和 $!A \ident (!B \lif !C)$ 的情形，证明 和~(b)。
\end{prob}

\begin{prob}
  证明每个经典重言式都是~$\LogHal$ 中的重言式。
\end{prob}

尽管它们与经典逻辑具有相同的重言式，其后承关系却不同。例如，$\LogLP$ 是\emph{次协调的}，因为 $\lnot p, p \Entails/ q$，所以爆炸原理 $\lnot !A, !A \Entails !B$ 并非普遍成立。（它对某些 $!A$ 和 $!B$ 是成立的，例如当 $!B$ 是重言式时。）

\begin{prob}
  下列后承关系在 (a)~$\LogLP$ 和 (b)~$\LogHal$ 中哪些成立？请为每一项构造真值表。
  \begin{enumerate}
    \item $p, p \lif q \Entails q$
    \item $\lnot q, p \lif q \Entails \lnot p$
    \item $p \lor q, \lnot p \Entails q$
    \item $\lnot p, p \Entails q$
    \item $p \Entails p \lor q$
    \item $p \lif q, q\lif r \Entails p \lif r$
  \end{enumerate}
\end{prob}

如果在 $\LogLuk[3]$ 中把 $\Undef$ 也指定为指定值，会怎样？

\begin{defn}
  \emph{三值 R-Mingle 逻辑}~$\LogRM[3]$ 由以下矩阵定义：
  \begin{enumerate}
    \item 标准命题语言 $\Lang L_0$，包含联结词
    $\lfalse$、$\lnot$、$\land$、$\lor$、$\lif$。
    \item 真值集合为 $V = \{\True, \Undef, \False\}$。
    \item $\True$ 和 $\Undef$ 是指定值，即 $V^+ = \{\True, \Undef\}$。
    \item 真值函数与 Łukasiewicz 逻辑~$\LogLuk[3]$ 中的相同。
  \end{enumerate}
\end{defn}

\begin{prob}
  下列后承关系在 $\LogRM[3]$ 中哪些成立？
  \begin{enumerate}
    \item $p, p \lif q \Entails q$
    \item $p \lor q, \lnot p \Entails q$
    \item $\lnot p, p \Entails q$
    \item $p \Entails p \lor q$
  \end{enumerate}
\end{prob}

有时，仅仅通过改变指定值，不同的真值表就能生成相同的逻辑（语义后承关系）。
例如，如果在 Gödel 逻辑中取 $V^+ = \{\True,
\Undef\}$ 而不是~$\{\True\}$ 作为指定值集合，就会出现这种情况。

\begin{prop}\ollabel{prop:gl-udes}
  由 $V = \{\False, \Undef,\True\}$、$V^+=\{\True,
  \Undef\}$ 及三值 Gödel 逻辑的真值函数构成的矩阵，定义的是经典逻辑。
\end{prop}

\begin{proof}
  练习。
\end{proof}

\begin{prob}
  通过证明以下结论来证明 \olref[mvl][thr][mul]{prop:gl-udes}：对于按 Gödel 逻辑方式定义但指定值为 $V^+=\{\True,\Undef\}$ 的逻辑~$\Log L$，若 $\Gamma \Entails/[\Log L] !B$ 则 $\Gamma \Entails/[\LogCL] !B$。
  请运用 \olref[mvl][thr][mul]{prop:LP-taut-CL} 中的思路，但无需证明性质 和~(b)，而是证明 $\pValue v(!A)[\LogGod] = \False$ 当且仅当 $\pValue {v'}(!A)[\LogCL] = \False$（从而 $\pValue
  v(!A)[\LogGod] \in \{\True,\Undef\}$ 当且仅当 $\pValue {v'}(!A)[\LogCL] =
  \True$），并解释为何这能确立该命题。
\end{prob}

\end{document}